\section{Certificate conditions and the global implication}\label{ms:certificates}

This section states the conditional global theorem, explains what a
finite certificate must establish, and records which conditions have been
verified.

\subsection{The mathematical interface}

Write $a_n=C_n-\CE$ and
$d_n=n(C_{n-1}-C_n)_+$.

\begin{assumption}[Finite base]\label{ms:base41}
For every integer $1\le j\le41$, one has $C_j\le\CE$.
This must be proved without using Assumption~\ref{ms:local42}.
\end{assumption}

\begin{assumption}[Equality-inclusive local exclusion]\label{ms:local42}
For every integer $n\ge42$, the following conditions do not hold
simultaneously:
\begin{equation}\label{ms:local-entry}
 \begin{gathered}
 0\le n(C_n-\CE)\le0.54,\qquad d_n=0,\qquad C_n<0.4226,\\
 C_{n-1}<0.4253,\qquad C_{n-2}<0.4253.
 \end{gathered}
\end{equation}
Since $C_n$ is attained (Lemma~\ref{ms:attainment}), this is equivalent to
the nonexistence of an extremal pair at index $n$ under these conditions,
which is how it would be proved.
\end{assumption}

A proof of Assumption~\ref{ms:local42} would combine the analytic results of
Sections~\ref{ms:support}--\ref{ms:counts} with finite certificates, in
four stages. First, the contact inequality and the moment
bounds give the initial support interval. Second, an ordered sequence of
Fourier estimates, moment intersections and support contractions gives
one closed parameter panel containing the maximizing law. Third, a joint
defect inequality and a confinement argument bound its conditional noise
and its support near the two-point skeleton. Fourth, the exact
conditional-count representation gives a comparison on every remaining
closed phase and noise cell. The analytic steps are in Sections~\ref{ms:support}--\ref{ms:counts}
and the appendices; their numerical inputs have not been independently
re-verified.

The stages are applied in order to one fixed extremal pair: each support
contraction may use only moment bounds proved from the preceding support
interval, and the deleted sums $S_{n-1}$ and $S_{n-2}$ are compared with
the Gaussian through the constants $C_{n-1}$ and $C_{n-2}$.

\begin{theorem}[Conditional all-index bound]\label{ms:conditional-global}
Suppose Assumptions~\ref{ms:base41} and~\ref{ms:local42} hold. Then
\[
 C_n<\CE\quad(n\ge42),\qquad C^*=\CE.
\]

\end{theorem}
\begin{proof}
If the strict tail conclusion fails, choose the least $n\ge42$ with
$C_n\ge\CE$. Both predecessors are at most $\CE$: at $n=42$ these are
exactly the base indices $40$ and $41$. Hence $d_n=0$ and the two
deleted caps in~\eqref{ms:local-entry} hold. The smoothing recurrence
of Section~\ref{ms:smoothing} gives
\[
 0\le n(C_n-\CE)<\frac{3003}{10250}<0.3,
 \qquad C_n<\frac{73}{175}<0.4226.
\]
Thus all conditions in~\eqref{ms:local-entry} hold, contradicting
Assumption~\ref{ms:local42}. The base
and tail give $C^*\le\CE$; the fixed Bernoulli limit
\eqref{ms:bernoulli-limit} gives the reverse inequality.
\end{proof}

If Assumption~\ref{ms:local42} is weakened to exclude only $C_n>\CE$, the
same argument, applied to the least $n$ with $C_n>\CE$, still gives
$C^*=\CE$.

\subsection{What a finite certificate must establish}

A certificate has three ingredients. Its domain
description contains every actual parameter tuple allowed by the
analytic hypotheses. Its local proof establishes the required inequality
throughout each closed cell. Its coverage proof shows that the cells
exhaust that domain.

The following elementary lemma is the basic tool for tensor certificates.

\begin{lemma}[Common-marginal expectation and score]
\label{ms:certificate-expectation}
Let $K\subset\R$ contain the support of $P\in\mathcal P$. Let
$f_0,\ldots,f_d$ be integrable functions under $P$, with $f_0=1$, and let
$m_j=\E f_j(X)$. For a finite tensor sum
\[
 T(x_1,\ldots,x_n)=\sum_{\alpha}t_\alpha
                              \prod_{i=1}^n f_{\alpha_i}(x_i),
\]
independence gives
\begin{equation}
 \E T(X_1,\ldots,X_n)=\sum_\alpha t_\alpha\prod_{i=1}^n m_{\alpha_i}.
\label{ms:tensor-expectation}
\end{equation}
If $T\ge\mathbf1_{\{x_1+\cdots+x_n\le H\}}$ on all of $K^n$ and
\[
 \sum_\alpha t_\alpha\prod_i m_{\alpha_i}
         -\Phi(H/\sqrt n)\le a\rho/\sqrt n,
\]
then $\Pp(S_n\le H)-\Phi(H/\sqrt n)\le a\rho/\sqrt n$.
\end{lemma}
\begin{proof}
Each product is absolutely integrable, because independence factors its
absolute expectation into the finite quantities $\E|f_j(X)|$.
Fubini and linearity prove \eqref{ms:tensor-expectation}. Integrate the
pointwise majorant and use $\rho>0$ to obtain the score inequality.

\end{proof}

For example, a tensor product containing a centered factor has zero
expectation when all other factors concern different independent variables.
Such terms may improve a pointwise majorant without altering its expectation.
Polynomial factors of degree above three require the corresponding moments.
For the extremal laws used here, these are finite because support is compact.

An exact polynomial positivity certificate often uses Bernstein coefficients.
On $[0,1]$, set $B_{j,d}(t)=\binom dj t^j(1-t)^{d-j}$. They are
nonnegative and sum to one. The identity
\[
 t^k=\sum_{j=k}^d\frac{\binom jk}{\binom dk}B_{j,d}(t)
\]
follows from
$\binom dj\binom jk=\binom dk\binom{d-k}{j-k}$ and the binomial theorem.
Applying it in each coordinate proves that a polynomial
$\sum_\alpha a_\alpha t^\alpha$ of coordinate degrees at most $d_i$ has
tensor Bernstein coefficients
\[
 b_\nu=\sum_{\alpha\le\nu}a_\alpha
                      \prod_i\frac{\binom{\nu_i}{\alpha_i}}
                                      {\binom{d_i}{\alpha_i}}.
\]
Thus nonnegative coefficients prove nonnegativity on the whole closed box.
For $P(x)=\sum_{i=0}^d a_i x^i$ on $[l,h]$, the change
$x=l+(h-l)t$ gives power coefficients
$\widetilde a_j=(h-l)^j\sum_{i=j}^d a_i\binom ij l^{i-j}$.
Applying the preceding identity gives the Bernstein coefficients on that
interval. Rational affine changes and subdivision preserve exactness. A certificate must also verify the polynomial identity and the coverage
of the domain.

The event plane needs particular care. On the closed side
$\sum x_i\le H$ the required inequality is $T-1\ge0$; on the other side
it is $T\ge0$. A certificate which verifies only $T\ge0$ at a boundary
point has not proved the closed-event majorant there. Sign orthants,
event-plane intersections and endpoint atoms must be included in the
pointwise proof.

\begin{proposition}[Closed-cover soundness]
\label{ms:closed-cover}
Suppose a family of actual laws and thresholds has an analytically proved
map into a covered parameter domain. A finite certificate has finitely
many roots covering that domain and a finite subdivision tree at each
root. Every split retains both closed children. At each leaf, either
a proved necessary condition excludes every actual mapped point in the
leaf, or a valid whole-leaf inequality bounds its positive score by $a$.
Then all actual points in the root domain have positive score at most $a$.
If every leaf instead excludes a putative attained score $C\ge\CE$, no
such attained witness belongs to the root domain.
\end{proposition}
\begin{proof}
An actual point at a parent belongs to at least one child, including if it
lies on the common boundary. Induction up each finite tree transfers the
leaf conclusion to its root. The root union transfers it to the complete
mapped domain. In the second version no actual point can lie in any leaf, a
contradiction.
\end{proof}

A root may describe a relaxation containing parameters of no actual law;
this does not affect soundness, whereas a missing root or an unfinished
subdivision does. A minimum of convex bounds need not be convex, so an
endpoint maximization must be justified for each bound before the minimum
is taken.

In interval computations every transcendental function is enclosed
outwards and every denominator is proved positive on the whole cell; a
lower bound for a binomial mass multiplies a bracket only after the
bracket is known to be negative.

A cover must include both signs of the phase, all admissible nearest
counts, the endpoints of probability panels and the boundaries between
noise bands. The case $\lambda=n\E U^2=0$ is separate: then $U=0$ almost
surely, the law is a two-point law, and Schulz's theorem
applies~\cite{Schulz2016}.

\subsection{Status of the remaining conditions}

Sections~\ref{ms:extremizers}--\ref{ms:counts} prove the analytic
implications stated there, with their moment, support and remainder
hypotheses. Bounded exact checks accompany the support, curvature and
score-excess calculations. The panel consequences in
Section~\ref{ms:resources} remain conditional on the coefficient and
moment bounds from Section~\ref{ms:fourier:rescore:re-atlas}.

The finite base (Assumption~\ref{ms:base41}) has not been established
in this work. The appendices give reductions for individual indices,
but the complete independent proof through $41$ is missing. An older
assembly for $40\le n\le42$ uses tail and coarse-ceiling inputs; the
current independent assembly named by the source material was unavailable.

The local exclusion (Assumption~\ref{ms:local42}) also remains open.
Some inherited local statements begin at $47$ or have finite upper limits.
Appendix~\ref{ms:uniform-bridge} gives the parameter changes needed at
$42$. Completing the exclusion requires the corresponding analytic bounds
and a complete finite cover at that scope. The longer Fourier,
confinement and finite-index arguments are included for review; their
whole chain and numerical inputs have not been independently verified
in full here.

\begin{center}
\begin{tabular}{@{}P{0.36\textwidth}P{0.30\textwidth}P{0.27\textwidth}@{}}
\toprule
Result & Mathematical status & Computation status\\
\midrule
Attainment, contact and smoothing & Proved for the stated class of laws &
Bounded checks recorded\\[3pt]
Support, conserved identity, score charge and curvature & Proved under the
stated hypotheses & Exact algebra and bounded arithmetic checked\\[3pt]
Complete Fourier, confinement and finite-index chain & Supplied arguments;
full independent review remains & No complete replay here\\[3pt]
Finite base and uniform local exclusion & Open & No complete acceptance\\
\bottomrule
\end{tabular}
\end{center}
